\chapter[Korelasyon, Regresyon ve Uygulama]{Korelasyon ve Basit Doğrusal Regresyon; Nedensellik Ayrımı ve Teknoloji Destekli Uygulama}
\label{ch:correlation-regression}

\begin{prismbox}[PrismLight]{Öğrenme çıktıları}
Bu bölümün sonunda:
\begin{itemize}
  \item iki nicel değişken arasındaki yönü, biçimi, gücü ve olağandışı gözlemleri saçılım grafiğinde değerlendirebilecek,
  \item Pearson korelasyonunu, regresyon eğimini, $R^2$ ve güven aralıklarını hesaplayıp yorumlayabilecek,
  \item artık grafiği üzerinden doğrusallık, sabit varyans ve etkili gözlem sorunlarını tanıyabilecek,
  \item ilişki, öngörü ve nedensel etki iddialarını araştırma tasarımına göre birbirinden ayırabileceksiniz.
\end{itemize}
\end{prismbox}

\begin{prerequisitebox}
İki nicel değişken, ortalama, standart sapma, standartlaştırma, güven aralığı
ve hipotez testi kavramları bilinmelidir. Koordinat sistemi ve doğru denklemi
hatırlanmalıdır.
\end{prerequisitebox}

İki nicel değişken için grafik seçimi Bölüm~\ref{sec:descriptive-graphs}'te,
güven aralığı ve test yorumu Bölüm~\ref{ch:confidence-intervals} ile
Bölüm~\ref{ch:hypothesis}'te hazırlanmıştır.

\section{Saçılım grafiği ve korelasyon}

İki nicel değişken arasındaki ilişki önce saçılım grafiğiyle incelenir.
Pearson örneklem korelasyonu
\[
r=\frac{\sum_i(x_i-\bar x)(y_i-\bar y)}
{\sqrt{\sum_i(x_i-\bar x)^2\sum_i(y_i-\bar y)^2}}
\]
olup $-1$ ile $1$ arasındadır. İşaret ilişkinin yönünü, mutlak değer doğrusal
ilişkinin gücünü gösterir.

Korelasyon birimsizdir ve $X$ ile $Y$ yer değiştirdiğinde değişmez. Ölçüm
birimlerinin metre yerine santimetreye çevrilmesi de $r$ değerini değiştirmez.
Buna karşılık regresyon eğimi değişkenlerin birimlerine bağlıdır ve $Y$'nin
$X$'e göre regresyonuyla $X$'in $Y$'ye göre regresyonu aynı doğru değildir.

\section{Korelasyon ve nedensellik}
\label{sec:correlation-causality}

Korelasyon tek başına nedensel etki göstermez. Ters nedensellik, ortak neden,
seçim yanlılığı veya tesadüf gözlenen ilişkiyi açıklayabilir. Nedensel yorum
araştırma tasarımı, zaman sırası ve karıştırıcı değişkenler hakkında ek kanıt
gerektirir \citep{shadish2002,gelman2020}.

\begin{mistakebox}
$r=0{,}70$ değeri, bir değişkenin diğerinin yüzde 70'ini açıkladığı anlamına
gelmez. Basit doğrusal modelde açıklanan varyans oranı $R^2=r^2=0{,}49$ olur.
\end{mistakebox}

\section{Basit doğrusal regresyon}
\label{sec:simple-regression}

Basit doğrusal regresyon
\[
Y_i=\beta_0+\beta_1X_i+\varepsilon_i
\]
modeliyle $X$ ve $Y$ arasındaki ortalama doğrusal ilişkiyi açıklar. $\beta_1$,
$X$ bir birim arttığında beklenen $Y$ değerindeki ortalama değişimdir.
Uydurulan doğru $\hat Y=b_0+b_1X$ biçimindedir.

Basit doğrusal regresyonda eğim ve sabit terim
\[
b_1=r\frac{s_Y}{s_X},
\qquad
b_0=\bar y-b_1\bar x
\]
eşitlikleriyle ilişkilidir. Eğim, $X$'te bir birimlik farkla ilişkili ortalama
$Y$ farkını gösterir. Sabit terim ise $X=0$ için modelin ortalama yanıtıdır;
sıfır değeri gözlenen aralıkta değilse çoğu zaman doğrudan bağlamsal anlam
taşımaz.

Basit doğrusal regresyonda $H_0:\beta_1=0$ testi ile
$H_0:\rho=0$ Pearson korelasyon testi aynı iki yönlü \pval-değerini verir:
\[
t=r\sqrt{\frac{n-2}{1-r^2}},\qquad sd=n-2.
\]
Bu eşdeğerlik yalnız aynı gözlemler ve doğrusal model için geçerlidir.

\section{Uyum ve model kontrolü}
\label{sec:regression-diagnostics}

$R^2$, yanıt değişkenliğinin model tarafından açıklanan oranını özetler.
Yüksek $R^2$ nedensellik veya doğru model garantisi değildir. Doğrusallık,
sabit hata varyansı, bağımsızlık, aykırı değerler ve artıkların davranışı
grafiklerle değerlendirilmelidir.

Artık $e_i=y_i-\hat y_i$ olarak tanımlanır. Artık--uydurulan değer grafiğinde
sıfır çevresindeki rastgele bant doğrusal ve sabit varyanslı yapıyı destekler.
Eğri örüntü doğrusal olmayan ilişkiye, giderek açılan huni değişen hata
varyansına, tek başına uzak bir nokta ise aykırı gözleme işaret edebilir.
Artıkların normal dağılımı özellikle küçük örneklemde katsayı testleri ve
aralıklar için önemlidir; doğrusallığı değerlendirmek için yalnız normal
olasılık grafiğine bakmak yeterli değildir.

Uç bir $X$ değerine sahip gözlem yüksek kaldıraçlıdır. Böyle bir gözlem aynı
zamanda regresyon çizgisinden uzaksa eğimi ve diğer katsayıları belirgin
biçimde değiştirebilir. Cook uzaklığı gibi ölçüler etkili gözlemleri tanımaya
yardım eder; ancak yüksek değer gözlemin yanlış olduğunu kanıtlamaz.

\section{Gerçek eğitim verisi üzerinde uygulama}

On altı öğrenciden oluşan öğretim amaçlı veri setinde haftalık çalışma süresi
ve son test puanı incelendiğinde
$r(14)=0{,}981$, $p<0{,}001$ bulunmuştur. Uydurulan regresyon doğrusu
\[
\widehat{\text{Son puan}}=44{,}60+3{,}14\times\text{Saat}
\]
biçimindedir. Gözlenen aralıkta bir saatlik çalışma süresi artışı, ortalama
son test puanında yaklaşık $3{,}14$ puanlık artışla ilişkilidir. Gözlemsel
veri nedeniyle bu katsayı nedensel etki olarak yorumlanmaz.
Bu öğretim verisinin kaynağı ve kullanım bağlamı için bkz. \citet{gill2001}.

\begin{figure}[H]
\centering
\begin{tikzpicture}[x=1.12cm,y=0.22cm]
  \draw[->,PrismGray] (1.5,0)--(8.7,0)
    node[right,font=\scriptsize]{Çalışma süresi};
  \draw[->,PrismGray] (1.5,0)--(1.5,23)
    node[above,font=\scriptsize]{Son puan};
  \foreach \x in {2,...,8}
    \node[below,font=\scriptsize] at (\x,0) {\x};
  \foreach \y/\label in {2/50,7/55,12/60,17/65,22/70} {
    \draw[PrismGray] (1.43,\y)--(1.57,\y);
    \node[left,font=\scriptsize] at (1.43,\y) {\label};
  }
  \foreach \x/\y in {5/12,3/6,7/20,6/15,6/17,2/4,8/22,4/10,
                        4/9,5/13,3/4,7/16,3/7,5/12,2/2,6/15}
    \fill[PrismSubheadingBlue!85!black] (\x,\y) circle (2.1pt);
  \draw[PrismBlue,very thick] (2,2.87)--(8,21.70);
\end{tikzpicture}
\caption{Çalışma süresi, son test puanı ve uydurulan regresyon doğrusu.}
\label{fig:simple-regression-scatter}
\end{figure}

Şekil~\ref{fig:simple-regression-scatter}'deki doğrusal örüntü için modelin
$R^2$ değeri $0{,}962$'dir. Bu, gözlenen örneklemde son test
puanlarındaki değişkenliğin yaklaşık yüzde 96,2'sinin çalışma süresiyle
kurulan doğrusal model tarafından açıklandığını gösterir. Küçük ve seçilmiş
bir örneklemdeki yüksek değer, yeni öğrencilerde aynı başarının elde
edileceğinin garantisi değildir. Sabit terim 44,60 puandır; ancak sıfır saat
veri aralığının dışındaysa bu değer bağımsız bir eğitimsel sonuç olarak
yorumlanmamalıdır.

Eğim tahmininin standart hatası 0,168 ve yüzde 95 güven aralığı yaklaşık
$[2{,}78;3{,}50]$ puandır. Bu aralık, gözlenen veri ve model koşulları altında
bir saatlik çalışma süresi artışıyla ilişkili ortalama puan değişimi için
makul değerleri gösterir. Korelasyonun yüzde 95 güven aralığı da yaklaşık
$[0{,}944;0{,}993]$'tür; küçük örnekleme rağmen doğrusal ilişkinin güçlü
olduğuna işaret eder.

\section{Çözümlü öngörü ve artık hesabı}

Altı saat çalışan bir öğrenci için modelin öngördüğü puan
\[
\hat y=44{,}60+3{,}14(6)=63{,}44
\]
olarak bulunur. Veri setinde altı saat çalışan ve 65 puan alan bir öğrenci
için artık
\[
e=y-\hat y=65-63{,}44=1{,}56
\]
puandır. Pozitif artık, öğrencinin modelin öngördüğünden daha yüksek puan
aldığını gösterir. Yuvarlanmamış katsayılarla hesaplandığında değerler
sırasıyla 63,42 ve 1,58'dir.

Tek bir öğrenci için 63,44 puan ``kesin sonuç'' değildir. Regresyon doğrusu
ortalama ilişkiyi verir; bireysel öngörüde artık değişkenliği de bulunduğu
için öngörü aralığı kullanılmalıdır.

\begin{figure}[H]
\centering
\begin{tikzpicture}[x=0.38cm,y=0.85cm]
  \draw[->,PrismGray] (0,-3.2)--(0,3.2) node[above,font=\scriptsize]{artık};
  \draw[->,PrismGray] (0,0)--(22.5,0)
    node[above left,font=\scriptsize]{uydurulan değer};
  \foreach \x/\label in {1/50,6/55,11/60,16/65,21/70} {
    \draw[PrismGray] (\x,-0.08)--(\x,0.08);
    \node[below,font=\tiny] at (\x,-0.08) {\label};
  }
  \foreach \y in {-3,-2,-1,1,2,3} {
    \draw[PrismGray] (-0.18,\y)--(0.18,\y);
    \node[left,font=\tiny] at (-0.18,\y) {\y};
  }
  \draw[PrismWarnOrange!85!black,dashed] (0,0)--(22,0);
  \foreach \x/\y in {11.28/-.28,5.01/-.01,17.56/1.44,14.42/-.42,
    14.42/1.58,1.87/1.13,20.70/.30,8.15/.85,8.15/-.15,11.28/.72,
    5.01/-2.01,17.56/-2.56,5.01/.99,11.28/-.28,1.87/-.87,14.42/-.42}
    \fill[PrismSubheadingBlue!85!black] (\x,\y) circle (2.0pt);
\end{tikzpicture}
\caption{Uydurulan değerlere karşı artıklar.}
\label{fig:simple-regression-residuals}
\end{figure}

Şekil~\ref{fig:simple-regression-residuals}'de artıkların sıfır çevresinde
belirgin bir eğri veya huni oluşturmadan
dağılması doğrusal biçim ve sabit varyans varsayımlarını destekler.

\section{Yazılım uygulaması: korelasyon, regresyon ve tanı grafikleri}

\begin{softwarebox}
\begin{lstlisting}
import numpy as np
import matplotlib.pyplot as plt
from scipy import stats

hours = np.array([5,3,7,6,6,2,8,4,4,5,3,7,3,5,2,6])
post = np.array([60,54,68,63,65,52,70,58,
                 57,61,52,64,55,60,50,63])

print(stats.pearsonr(hours, post))
fit = stats.linregress(hours, post)
print(fit)
print("R-kare:", fit.rvalue**2)
print("6 saat icin ongoru:", fit.intercept + fit.slope * 6)

fitted = fit.intercept + fit.slope * hours
residuals = post - fitted

fig, axes = plt.subplots(1, 2, figsize=(8, 3))
axes[0].scatter(hours, post)
axes[0].plot(hours, fitted, color="#1F4E79")
axes[0].set(xlabel="Calisma suresi", ylabel="Son test puani")
axes[1].scatter(fitted, residuals)
axes[1].axhline(0, color="#1F4E79", linestyle="--")
axes[1].set(xlabel="Uydurulan deger", ylabel="Artik")
plt.tight_layout()
plt.show()
\end{lstlisting}
SPSS'te önce \emph{Graphs $>$ Chart Builder} ile saçılım grafiği, ardından
\emph{Analyze $>$ Correlate $>$ Bivariate} ve
\emph{Analyze $>$ Regression $>$ Linear} yolları izlenebilir. R'de
\texttt{cor.test(hours, post)} ve \texttt{lm(post \textasciitilde{} hours)}
aynı temel soruları yanıtlar. Her ortamda artık grafikleri ayrıca
incelenmelidir. Python hesaplarının dayandığı temel bilimsel yazılım
altyapıları için bkz. \citet{harris2020} ve \citet{virtanen2020}.
\end{softwarebox}

\begin{assumptionbox}
Korelasyon ve eğim aykırı ya da yüksek etkili bir gözlemden güçlü biçimde
etkilenebilir. Saçılım grafiği incelenmeden yalnız katsayıya bakılmamalıdır.
Doğrusal olmayan belirgin bir ilişki varsa $r$ sıfıra yakın olsa bile iki
değişken arasında güçlü bir örüntü bulunabilir.
\end{assumptionbox}

\section{Öngörü ve dışa taşırma}
\label{sec:extrapolation}

Regresyon doğrusu, veri aralığı içindeki yeni bir $X$ değeri için ortalama
yanıtı veya tek bir gözlemi öngörmede kullanılabilir. Tek gözlem için öngörü
aralığı, ortalama yanıtın güven aralığından daha geniştir. Gözlenen çalışma
süresi aralığının çok dışına taşırma yapmak, doğrusal ilişkinin o bölgede de
sürdüğü varsayımına dayanır ve çoğu zaman güvenilir değildir.

\section{Raporlama}

İyi bir rapor araştırma sorusunu, veri kaynağını, saçılım grafiğini,
korelasyon katsayısını, regresyon eğimini, güven aralığını, \pval-değerini,
$R^2$ değerini, varsayım kontrollerini ve nedensellik sınırlılığını birlikte
verir.

\section{Uygulama kontrol listesi}

\begin{enumerate}
  \item Her iki değişkenin nicel ve ölçümlerin bağımsız olduğunu doğrulayın.
  \item Saçılım grafiğinde yönü, biçimi, gücü ve aykırı değerleri inceleyin.
  \item Korelasyon katsayısını güven aralığı ve örneklem hacmiyle raporlayın.
  \item Regresyon eğimini değişkenlerin birimleriyle yorumlayın.
  \item Artık grafiğinde eğrilik, değişen yayılım ve etkili gözlemleri araştırın.
  \item Nedensel yorumun araştırma tasarımı tarafından desteklenip desteklenmediğini belirtin.
\end{enumerate}

\begin{outputguidebox}
Regresyon çıktısında ``Model Summary'' bölümündeki $R^2$ ile katsayılar
tablosundaki eğimi birbirine karıştırmayın. Eğim, $X$'teki bir birim değişimle
ilişkili $Y$ değişimini; $R^2$ ise örneklemde açıklanan değişkenlik oranını
gösterir. Katsayının standart hatası, güven aralığı ve \pval-değeri aynı satırda
birlikte okunmalıdır.
\end{outputguidebox}

\section{Çözümlü problemler}

\subsection*{Problem 1: Saçılım grafiğini sistematik okuma}

Bir saçılım grafiğinde çalışma süresi arttıkça başarı puanları genel olarak
artmakta; noktalar yaklaşık bir doğru çevresinde toplanmakta, ancak en yüksek
çalışma süresine sahip bir öğrenci çizginin oldukça altında kalmaktadır.
Grafiği yön, biçim, güç ve olağandışı gözlem bakımından yorumlayın.

\textbf{Çözüm.} Genel eğilim pozitif yönlüdür: yüksek çalışma süreleri daha
yüksek puanlarla birliktedir. Noktaların doğru çevresinde toplanması ilişkinin
yaklaşık doğrusal olduğunu, çizgi çevresindeki saçılmanın sınırlı olması ise
doğrusal ilişkinin görece güçlü olabileceğini düşündürür. Bununla birlikte en
yüksek $X$ değerindeki öğrenci hem yüksek kaldıraçlı hem de büyük negatif
artıklı olabilir. Bu gözlem eğimi ve korelasyonu belirgin etkileyebilir.

Gözlem yalnız sıra dışı olduğu için silinmez. Veri giriş doğruluğu, öğrencinin
araştırma evrenine ait olup olmadığı ve ölçüm koşulları incelenir. Analiz
gözlem dâhil ve hariç yinelenerek duyarlılık değerlendirilir; iki sonuç
belirgin farklıysa her ikisi ve gerekçesi raporlanır.

\subsection*{Problem 2: Özetlerden regresyon doğrusunu kurma}

Bir araştırmada $\bar x=5$, $s_X=2$, $\bar y=70$, $s_Y=10$ ve $r=0{,}60$
bulunmuştur. $Y$'nin $X$'e göre regresyon doğrusunu kurun, $X=7$ için öngörü
yapın ve gözlenen $Y=79$ ise artığı hesaplayın.

\textbf{Çözüm.} Eğim
\[
b_1=r\frac{s_Y}{s_X}=0{,}60\frac{10}{2}=3
\]
ve sabit terim
\[
b_0=\bar y-b_1\bar x=70-3(5)=55
\]
olur. Uydurulan doğru
\[
\hat Y=55+3X
\]
biçimindedir. $X=7$ için öngörü $\hat y=55+3(7)=76$'dır. Gözlenen değer 79
ise artık
\[
e=y-\hat y=79-76=3
\]
olur; öğrenci modelin ortalama öngörüsünden üç puan daha yüksek sonuç
almıştır. Modelin açıklanan varyans oranı $R^2=0{,}60^2=0{,}36$'dır.

\subsection*{Problem 3: Korelasyon testi ve güven aralığı}

Yirmi beş öğrencide iki nicel değişken arasındaki Pearson korelasyonu
$r=0{,}40$ bulunmuştur. Sıfır korelasyon hipotezini sınayın ve Fisher
dönüşümüyle elde edilmiş yüzde 95 güven aralığı $[0{,}006;0{,}687]$ ise
sonucu yorumlayın.

\textbf{Çözüm.} Test istatistiği
\[
t=0{,}40\sqrt{\frac{25-2}{1-0{,}40^2}}
=0{,}40\sqrt{\frac{23}{0{,}84}}\approx2{,}09
\]
ve serbestlik derecesi 23'tür. Çift yönlü \pval-değeri yaklaşık 0,048'dir.
Yüzde 5 düzeyinde sıfır doğrusal korelasyona karşı sınırda kanıt vardır.

Güven aralığı sıfırın hemen üzerindeki çok zayıf ilişkilerden yaklaşık 0,69
düzeyindeki güçlü ilişkilere kadar geniş bir değer kümesini kapsamaktadır.
Dolayısıyla yalnız ``anlamlı korelasyon'' demek tahmin belirsizliğini gizler.
Korelasyonun pozitif görünmesi nedensel etkiyi kanıtlamaz.

\subsection*{Problem 4: Etkili gözlemin korelasyon ve eğime etkisi}

$X=(1,2,3,4,5)$ ve $Y=(2,4,6,8,20)$ verileri için yazılım
$r=0{,}894$ ve $\hat Y=-4+4X$ sonuçlarını vermektedir. Son gözlem olmadan ilk
dört nokta için $r=1$ ve $\hat Y=2X$ elde edilmektedir. Sonuçları açıklayın.

\textbf{Çözüm.} İlk dört gözlem kusursuz biçimde $Y=2X$ doğrusu üzerindedir.
$(5,20)$ noktası $X$ bakımından veri aralığının ucunda ve bu örüntünün
oldukça üzerindedir. Bu nedenle hem yüksek kaldıraçlı hem de etkili bir
gözlemdir. Tek bir nokta eğimi 2'den 4'e, sabit terimi 0'dan $-4$'e
değiştirmiş ve korelasyonu 1'den 0,894'e düşürmüştür.

Her iki korelasyon da güçlü görünmesine rağmen öngörü denklemleri önemli
ölçüde farklıdır. Bu örnek, katsayıların saçılım grafiği ve etki tanıları
olmadan yorumlanmaması gerektiğini gösterir. Son gözlemin doğruluğu
incelenmeli; geçerliyse keyfî olarak çıkarılmamalıdır.

\subsection*{Problem 5: Artık grafiğinde sorun tanıma}

Üç ayrı artık grafiğinde sırasıyla şu örüntüler gözlenmiştir: (a) belirgin U
biçimi, (b) uydurulan değer büyüdükçe açılan huni ve (c) sıfır çevresinde
rastgele bant içinde tek bir çok uzak nokta. Her örüntünün olası anlamını ve
ilk inceleme adımını yazın.

\textbf{Çözüm.}

\begin{enumerate}[label=\alph*)]
  \item U biçimi, ortalama ilişkinin doğrusal olmadığını düşündürür. İlk olarak
  özgün saçılım grafiği incelenmeli; kuramsal olarak anlamlıysa eğrisel terim
  veya uygun dönüşüm değerlendirilmelidir.
  \item Huni biçimi hata varyansının $X$ ya da uydurulan değer boyunca sabit
  olmadığını gösterir. Ölçek, alt gruplar ve veri toplama süreci incelenmeli;
  sağlam standart hatalar veya uygun modelleme seçenekleri düşünülmelidir.
  \item Tek uzak artık aykırı bir $Y$ gözlemine işaret eder. Kaldıraç ve Cook
  uzaklığıyla etkisi, kaynak kayıtla da ölçüm doğruluğu kontrol edilmelidir.
\end{enumerate}

Bu tanılar otomatik silme veya dönüşüm kuralları değildir. Model değişikliği
araştırma sorusu, ölçüm ölçeği ve veri üretim süreciyle gerekçelendirilmelidir.

\subsection*{Problem 6: Regresyon çıktısını bütüncül yorumlama}

Bir gözlemsel araştırmada çalışma süresinin başarı puanını açıklamasına yönelik
model için $n=80$, $b_1=1{,}80$, $\SE(b_1)=0{,}55$,
$t(78)=3{,}27$, $p=0{,}002$, yüzde 95 GA $[0{,}71;2{,}89]$ ve
$R^2=0{,}12$ verilmiştir. Artık grafiğinde hafif huni biçimi vardır. Sonucu
raporlayın.

\textbf{Çözüm.} Gözlenen veri aralığında çalışma süresindeki bir saatlik
artış, ortalama başarı puanında 1,80 puanlık artışla ilişkilidir. Güven
aralığı, model koşulları altında ilişkinin yaklaşık 0,71 ile 2,89 puan
arasında olabileceğini gösterir. Eğim sıfırdan farklıdır,
$t(78)=3{,}27$, $p=0{,}002$; ancak model başarı puanlarındaki örneklem
değişkenliğinin yalnız yüzde 12'sini açıklamaktadır.

Huni biçimi sabit varyans varsayımının zayıflayabileceğini gösterdiğinden
klasik standart hata ve aralıkların duyarlılığı incelenmelidir. Çalışma
gözlemsel olduğu için katsayı ``bir saat daha çalışmak başarıyı 1,80 puan
artırır'' biçiminde nedensel yorumlanamaz. Önceki başarı, motivasyon veya ders
desteği gibi karıştırıcılar ilişkiyi açıklayabilir.

\begin{outputguidebox}
Bir regresyon çıktısını okurken katsayının işaretini ve birimini, standart
hatasını, güven aralığını ve \pval-değerini aynı satırda değerlendirin. Ardından
$R^2$, örneklem hacmi, artık grafikleri ve etkili gözlem ölçülerini inceleyin.
Yazılımın katsayı üretmesi doğrusal biçimi, nedenselliği veya dışa taşırmanın
geçerliliğini kendiliğinden doğrulamaz.
\end{outputguidebox}

\section{Literatür verisiyle yazılım uygulamaları}
\label{sec:spss-b13}

\textbf{Amaç ve veri.} \texttt{b13.csv} ile birinci dönem notu
\texttt{g1} ve yıl sonu notu \texttt{g3} arasındaki doğrusal ilişkiyi
inceleyin \citep{cortez2008data}. Pearson korelasyonunda
$H_0:\rho=0$, sabit terimli basit regresyonda $H_0:\beta_1=0$
sınanır. Alternatifler iki yönlüdür; anlamlılık düzeyi 0,05,
katsayıların güven düzeyi yüzde 95'tir. Açıklayıcı değişken G1,
yanıt değişkeni G3'tür; bu yön bir nedensellik iddiası değildir.

\textbf{Ortak veri.} Karşılaştırmanın veri dosyası
\nolinkurl{companion/spss/csv/b13.csv} olup 395 satır ve 10 sütun
içerir. Kayıtların tamamı kullanılır; sıfır yıl sonu notu bulunan
38 kayıt eksik sayılmaz. B06'nın seçim göstergeleri ve ağırlığı
kullanılmaz. Sayısal uyum Excel içe aktarmanın doğrulandığı anlamına
gelmez. Kodları \texttt{companion} klasörünü içeren paket kökünden
çalıştırın; veri sözlüğü için bkz. sayfa~\pageref{sec:spss-guide}.

\subsection{IBM SPSS uygulaması}
\textbf{Başlamadan önce.} Aşağıdaki veri açma bloğunu
\texttt{EXECUTE} satırı dahil çalıştırın; sonra menü veya analiz
komutlarıyla devam edin. Filtre, ağırlık ve dosya bölme kapalı
olmalıdır. Dosya açılamazsa \texttt{/FILE} yolunu düzeltip yeniden
açın; önceki bölümden kalan etkin veriyle devam etmeyin.

\begin{spssadimlar}
\spssadim{Önce \emph{Graphs $\to$ Legacy Dialogs $\to$ Scatter/Dot} yolunda basit saçılım grafiği oluşturun. Yatay eksene \texttt{g1}, dikey eksene \texttt{g3} koyun; doğrusal örüntüyü ve aykırı kayıtları inceleyin.}
\spssadim{\emph{Analyze $\to$ Correlate $\to$ Bivariate} yolunu açın. \texttt{g1} ve \texttt{g3} değişkenlerini \emph{Variables} alanına taşıyın.}
\spssadim{Doğrusal ilişkiyi inceleyen bu uygulamada \emph{Pearson} ve \emph{Two-tailed} seçin. \emph{OK} ile tabloyu alın; yöntem seçimini çıkan $p$ değerine göre değiştirmeyin.}
\spssadim{Regresyon için \emph{Analyze $\to$ Regression $\to$ Linear} yolunu açın. \emph{Dependent} alanına \texttt{g3}, açıklayıcı değişken alanına \texttt{g1} koyun; yöntem \emph{Enter} olsun.}
\spssadim{\emph{Statistics} içinde \emph{Estimates}, \emph{Model fit} ve yüzde 95 \emph{Confidence intervals} seçin; \emph{Continue} seçin.}
\spssadim{\emph{Plots} içinde Y alanına \texttt{*ZRESID}, X alanına \texttt{*ZPRED} yerleştirin; standartlaştırılmış artık histogramını da seçin.}
\spssadim{\emph{Continue $\to$ OK} ile modeli çalıştırın. \emph{Coefficients} tablosunun standartlaştırılmamış B sütununu ve \emph{Model Summary} tablosunun $R^2$ değerini okuyun.}
\end{spssadimlar}

{\raggedright
\noindent\textbf{Komutların görevi.} \texttt{CORRELATIONS} Pearson ilişkisini, \texttt{REGRESSION} yıl sonu notunun doğrusal modelini hesaplar. \texttt{SCATTERPLOT} ve \texttt{RESIDUALS} model artıklarını incelemek içindir.\par}

\par\noindent\begin{minipage}{\linewidth}
\textbf{Uygulama kodu:} \texttt{b13}. Veri açma ve tanımlama:

\begin{lstlisting}[style=prismSPSS]
SET DECIMAL=DOT.
GET DATA /TYPE=TXT
 /FILE='companion/spss/csv/b13.csv'
 /ENCODING='UTF8'
 /ARRANGEMENT=DELIMITED /DELCASE=LINE
 /DELIMITERS="," /QUALIFIER='"'
 /FIRSTCASE=2
 /VARIABLES=id F8.0 school F8.0 sex F8.0 age F8.0
 studytime F8.0 absences F8.0 g1 F8.0 g2 F8.0
 g3 F8.0 pass10 F8.0.
FILTER OFF.
WEIGHT OFF.
SPLIT FILE OFF.
VARIABLE LABELS
 g1 'Ilk donem matematik notu'
 g3 'Yil sonu matematik notu'.
VARIABLE LEVEL
 id school sex pass10 (NOMINAL)
 studytime (ORDINAL)
 age absences g1 g2 g3 (SCALE).
EXECUTE.
\end{lstlisting}
\end{minipage}\par

\par\noindent\begin{minipage}{\linewidth}
\textbf{Korelasyon, model ve grafikler.} Veri açıldıktan sonra çalıştırın.

\begin{lstlisting}[style=prismSPSS]
DESCRIPTIVES VARIABLES=g1 g3
 /STATISTICS=MEAN STDDEV MIN MAX.
GRAPH /SCATTERPLOT(BIVAR)=g1 WITH g3.
CORRELATIONS /VARIABLES=g1 g3
 /PRINT=TWOTAIL /MISSING=PAIRWISE.
REGRESSION /STATISTICS=COEFF R ANOVA CI(95)
 /DEPENDENT=g3 /METHOD=ENTER g1
 /SCATTERPLOT=(*ZRESID,*ZPRED)
 /RESIDUALS=HISTOGRAM(ZRESID).
\end{lstlisting}

\end{minipage}\par

\begin{outputguidebox}
\textbf{Paylaşılan SPSS tablolarıyla doğrulanan değerler.}
\emph{Descriptive Statistics} tablosunda iki değişken için de
$n=395$ vardır. G1 ortalaması 10,91, standart sapması 3,319;
G3 ortalaması 10,42, standart sapması 4,581'dir.
\emph{Correlations} tablosunda $r=0{,}801$ ve iki yönlü
$p<0{,}001$ görülür. \emph{Model Summary}, $R^2=0{,}642$,
düzeltilmiş $R^2=0{,}641$ ve tahminin standart hatası 2,743 verir.

\emph{ANOVA} tablosunda $F(1,393)=705{,}842$ ve $p<0{,}001$;
\emph{Coefficients} tablosunda sabit $-1{,}653$, G1 eğimi
1,106 ve standartlaştırılmış Beta 0,801'dir. Eğim aralığı
$[1{,}024;\ 1{,}188]$ olarak gösterilir. B eğimi ile Beta
aynı sütun değildir. SPSS'in $p<0{,}001$ gösterimi sıfır demek
değildir; R ve Python daha küçük değerleri bilimsel gösterimle verir.

\emph{Residuals Statistics} tablosunda ham artık standart sapması
2,740, standartlaştırılmış artık standart sapması 0,999'dur.
Bunlar modelin 2,743 standart hatasıyla karıştırılmaz: artık
standart sapmasının böleni 394, model hata varyansının böleni
393'tür. Burada ZRESID, ham artığın model standart hatasına
bölünmesidir; standart sapmasının tam 1 olması beklenmez.
\end{outputguidebox}


\par\noindent\begin{minipage}{\linewidth}
\subsection{R uygulaması}
\label{sec:r-gercek-b13}

\texttt{lm} sabit terimli modeli kurar. Aşağıdaki blokları sırayla
çalıştırın; tahmin ve artıklar ham kayıtlardan ayrı nesnelerdir.

\begin{lstlisting}[style=prismR]
veri <- read.csv("companion/spss/csv/b13.csv")
stopifnot(nrow(veri) == 395, ncol(veri) == 10,
  !anyNA(veri), !anyDuplicated(veri$id),
  all(veri$pass10 == as.integer(veri$g3 >= 10)))
korelasyon <- cor.test(veri$g1, veri$g3,
  method="pearson", alternative="two.sided")
model <- lm(g3 ~ g1, data=veri)
ozet <- summary(model)
tahmin <- fitted(model)
artik <- residuals(model)
zpred <- as.numeric(scale(tahmin))
zresid <- artik/sigma(model)
print(c(satir=nrow(veri), sutun=ncol(veri),
  eksik=sum(is.na(veri)), sifir_not=sum(veri$g3 == 0)))
\end{lstlisting}
\end{minipage}\par

\textbf{Dosya yolu sorunu yaşarsanız.} İlk satır yerine aşağıdaki
komutla \texttt{b13.csv} dosyasını seçin; ardından kontrol bloğundan
devam edin.
\begin{lstlisting}[style=prismR]
veri <- read.csv(file.choose())
\end{lstlisting}

\par\noindent\begin{minipage}{\linewidth}
\textbf{Modelin sayısal özeti.} Küçük $p$ değerlerini sıfıra
yuvarlamamak için 15 anlamlı basamakla yazdırın.

\begin{lstlisting}[style=prismR]
print(c(r=unname(korelasyon$estimate),
  p_korelasyon=korelasyon$p.value,
  R2=ozet$r.squared, duzeltilmis_R2=ozet$adj.r.squared,
  model_sh=sigma(model),
  F=unname(ozet$fstatistic["value"]),
  df_model=unname(ozet$fstatistic["numdf"]),
  df_artik=df.residual(model),
  beta=coef(model)[["g1"]]*sd(veri$g1)/sd(veri$g3),
  G1_10_tahmin=unname(predict(model, data.frame(g1=10)))
), digits=15)
print(cbind(ozet$coefficients,
            confint(model, level=.95)), digits=15)
print(anova(model), digits=15)
\end{lstlisting}

\textbf{Çıktıyı okuma.} Paylaşılan R ekranlarında
$r=0{,}801467932017414$, $R^2=0{,}642350846052270$ ve
G1=10 için tahmin 9,40975711892797 bulunmuştur. Katsayı tablosu
hem sabit hem eğim için standart hata, $t$, $p$ ve güven aralığı
verir. R burada yeniden çalıştırılmamış; paylaşılan çıktılar incelenmiştir.
\end{minipage}\par

\par\noindent\begin{minipage}{\linewidth}
\textbf{Betimsel ve artık özetleri.} SPSS'in artık tablosuyla
karşılaştırma için aynı ölçeklemeyi koruyun.

\begin{lstlisting}[style=prismR]
betimsel <- function(degerler) {
  c(n=length(degerler), minimum=min(degerler),
    maksimum=max(degerler), ortalama=mean(degerler),
    ss=sd(degerler))
}
print(rbind(G1=betimsel(veri$g1), G3=betimsel(veri$g3)),
      digits=15)
print(rbind(tahmin=betimsel(tahmin), artik=betimsel(artik),
  zpred=betimsel(zpred), zresid=betimsel(zresid)), digits=15)
\end{lstlisting}

\textbf{Çıktıyı okuma.} R'de ZRESID aralığı yaklaşık
$[-4{,}236510;\ 1{,}828148]$, standart sapması 0,998730'dur.
Artık ortalamalarının $10^{-16}$ civarında görünmesi sayısal
hesap hassasiyetiyle ilgilidir; anlamlı bir ortalama sapması değildir.
\end{minipage}\par

\par\noindent\begin{minipage}{\linewidth}
\textbf{Grafikleri üretme.} Paylaşılan R görseli üç paneli de içerir.

\begin{lstlisting}[style=prismR]
par(mfrow=c(1, 3))
plot(veri$g1, veri$g3, xlab="G1", ylab="G3",
     main="G1 - G3", pch=16)
abline(model, col="blue", lwd=2)
plot(zpred, zresid, xlab="ZPRED", ylab="ZRESID",
     main="Artik - Tahmin", pch=16)
abline(h=0, lty=2)
hist(zresid, breaks=15, main="Artik histogrami",
     xlab="ZRESID")
par(mfrow=c(1, 1))
\end{lstlisting}
\end{minipage}\par

\par\noindent\begin{minipage}{\linewidth}
\subsection{Python uygulaması}
\label{sec:python-b13}

\texttt{linregress} eğim, sabit ve eğimin standart hatasını verir.
Eğim aralığında $n-2=393$ serbestlik derecesi kullanılır.
Blokları sırayla çalıştırın; \texttt{ddof=1} örneklem standart
sapmasını seçer.

\begin{lstlisting}[style=prismPythonApp]
import pandas as pd
import numpy as np
from scipy import stats

veri = pd.read_csv("companion/spss/csv/b13.csv")
assert veri.shape == (395, 10)
assert not veri.isna().any().any()
assert veri.id.is_unique
assert veri.pass10.eq((veri.g3 >= 10).astype(int)).all()
ilk_not = veri.g1.to_numpy()
son_not = veri.g3.to_numpy()
korelasyon = stats.pearsonr(ilk_not, son_not)
model = stats.linregress(ilk_not, son_not)
serbestlik = len(veri)-2
kritik = stats.t.ppf(.975, serbestlik)
egim_araligi = (model.slope
    + np.array([-1, 1])*kritik*model.stderr)
tahmin = model.intercept+model.slope*ilk_not
artik = son_not-tahmin
model_sh = np.sqrt(np.sum(artik**2)/serbestlik)
zresid = artik/model_sh
zpred = (tahmin-tahmin.mean())/tahmin.std(ddof=1)
\end{lstlisting}
\end{minipage}\par

\par\noindent\begin{minipage}{\linewidth}
\textbf{Sayısal çıktı.} Etiketler paylaşılan terminal çıktısıyla aynıdır.

\begin{lstlisting}[style=prismPythonApp]
print("Satir:", len(veri), "Sutun:", veri.shape[1])
print("Eksik:", int(veri.isna().sum().sum()))
print("Sifir not:", int((veri.g3 == 0).sum()))
ozet = {
    "Pearson r": korelasyon.statistic,
    "Korelasyon p": korelasyon.pvalue,
    "Sabit": model.intercept, "Egim": model.slope,
    "Egim standart hatasi": model.stderr,
    "Egim p": model.pvalue,
    "Egim GA alt": egim_araligi[0],
    "Egim GA ust": egim_araligi[1],
    "R2": model.rvalue**2,
    "Duzeltilmis R2": 1-(1-model.rvalue**2)
        * (len(veri)-1)/serbestlik,
    "Model standart hatasi": model_sh,
    "F": (model.slope/model.stderr)**2,
    "Artik serbestlik derecesi": serbestlik,
    "G1=10 ortalama tahmini": model.intercept+10*model.slope}
for ad, deger in ozet.items():
    print(f"{ad}: {deger:.15g}")
\end{lstlisting}

\textbf{Çıktıyı okuma.} Python ekranında 395 satır, 10 sütun,
sıfır eksik değer ve 38 sıfır not vardır. Korelasyon ve eğim için
$p\approx9{,}00143\times10^{-90}$ bulunmuştur. \texttt{.15g},
15 ondalık basamak değil 15 anlamlı basamak ister; çok küçük
değerleri bilimsel gösterimle korur.
\end{minipage}\par

\par\noindent\begin{minipage}{\linewidth}
\textbf{Grafikleri üretme.} Çizgi G1'e göre sıralanarak çizilir.

\begin{lstlisting}[style=prismPythonApp]
import matplotlib.pyplot as plt

sekil, eksenler = plt.subplots(1, 3, figsize=(15, 4))
sirali = np.argsort(ilk_not)
eksenler[0].scatter(ilk_not, son_not, alpha=.6)
eksenler[0].plot(ilk_not[sirali], tahmin[sirali],
                 color="red")
eksenler[0].set(xlabel="G1", ylabel="G3", title="G1 - G3")
eksenler[1].scatter(zpred, zresid, alpha=.6)
eksenler[1].axhline(0, color="black", linestyle="--")
eksenler[1].set(xlabel="ZPRED", ylabel="ZRESID",
                title="Artik - Tahmin")
eksenler[2].hist(zresid, bins=15, edgecolor="black")
eksenler[2].set(xlabel="ZRESID", ylabel="Frekans",
                title="Artik histogrami")
plt.tight_layout(); plt.show()
\end{lstlisting}
\end{minipage}\par

\Needspace{12\baselineskip}
\subsection{Sonuçların karşılaştırılması ve ortak rapor}
13 Eylül 2026'da paylaşılan SPSS tabloları ve grafikleri, R ekranları
ve Python terminal çıktısı ile üç panelli Python grafiği incelenmiştir.
Python'un 14 sayısal model özeti R'deki karşılıklarıyla uyumludur:
$p$ dışındaki değerlerde $10^{-10}$ mutlak, çok küçük $p$ değerlerinde
$10^{-10}$ bağıl tolerans kullanılmıştır. Bütün basamaklarda eşitlik
iddia edilmez. Python kodu proje CSV'siyle ayrıca çalıştırılmış;
SPSS sonuçları gösterilen yuvarlama hassasiyetlerinde uyumlu bulunmuştur.
Sabitin aralığı ve ayrıntılı ANOVA gibi Python terminalinde ayrıca
yazdırılmayan değerler için SPSS ve R çıktıları esas alınmıştır.
Aşağıdaki tablolar bu çıktıların kitap düzenine uyarlanmış özetidir;
ekran görüntüsü değildir.

\begin{table}[H]
\centering
\caption{B13'te korelasyon ve basit regresyonun uyumu}
\label{tab:b13-dogrulanmis-uyum}
\begin{tabular}{@{}lr@{}}
\toprule
Ölçü & Değer \\
\midrule
Gözlem sayısı & 395 \\
Pearson $r$ & 0,801468 \\
$R^2$ & 0,642351 \\
Düzeltilmiş $R^2$ & 0,641441 \\
Modelin artık standart hatası & 2,743359 \\
$F(1,393)$ & 705,842247 \\
Korelasyon ve eğim için $p$ (yaklaşık) & $9{,}00143\times10^{-90}$ \\
Standartlaştırılmış Beta & 0,801468 \\
G1=10 için ortalama tahmini & 9,409757 \\
\bottomrule
\end{tabular}
\par\smallskip
{\footnotesize Sabit terimli, tek açıklayıcılı bu modelde
$R^2=r^2$, standartlaştırılmış Beta $=r$ ve eğimin $t$ istatistiğinin
karesi modelin $F$ istatistiğidir. Korelasyon ve eğim testleri
bağımsız iki kanıt olarak sayılmaz.\par}
\end{table}

\begin{table}[H]
\centering
\caption{B13'te standartlaştırılmamış regresyon katsayıları}
\label{tab:b13-dogrulanmis-katsayi}
\begin{tabular}{@{}lrr@{}}
\toprule
Ölçü & Sabit & G1 eğimi \\
\midrule
B & $-1{,}652804$ & 1,106256 \\
Standart hata & 0,474745 & 0,041639 \\
$t(393)$ & $-3{,}481452$ & 26,567692 \\
İki yönlü $p$ (yaklaşık) & 0,000555 & $9{,}00143\times10^{-90}$ \\
\%95 güven aralığı: alt & $-2{,}586162$ & 1,024393 \\
\%95 güven aralığı: üst & $-0{,}719445$ & 1,188119 \\
\bottomrule
\end{tabular}
\par\smallskip
{\footnotesize Güven aralıkları klasik regresyon modeline dayanır;
dayanıklı standart hatalar kullanılmamıştır. Katsayıların yorumunda
artık grafikleri ve bağımsızlık varsayımı ayrıca değerlendirilir.\par}
\end{table}

\begin{table}[H]
\centering
\caption{B13'te regresyon varyans çözümlemesi}
\label{tab:b13-dogrulanmis-anova}
\begin{tabular}{@{}lrrr@{}}
\toprule
Kaynak & Kareler toplamı & sd & Kareler ortalaması \\
\midrule
Regresyon & 5312,182953 & 1 & 5312,182953 \\
Artık & 2957,725907 & 393 & 7,526020 \\
Toplam & 8269,908861 & 394 & --- \\
\bottomrule
\end{tabular}
\par\smallskip
{\footnotesize $F$, regresyon kareler ortalamasının artık kareler
ortalamasına oranıdır. Artık kareler ortalamasının karekökü 2,743359
model standart hatasını verir. Ara hesaplar yuvarlanmamış değerlerle
yapılır; tablolarda çoğu değer altı ondalık basamakla gösterilmiştir.\par}
\end{table}

\textbf{Grafik kontrolünün kapsamı.} Üç yazılımda da G1--G3
saçılımı, standartlaştırılmış artık--tahmin grafiği ve artık
histogramı paylaşılmıştır. Saçılımda pozitif doğrusal örüntü
görülür; G3=0 satırındaki noktalar ana örüntüden ayrılır.
Notlar tam sayı olduğundan noktalar üst üste binebilir; grafikteki
görünür simgelerden 395 kaydın veya sıfır notların sayısı okunmaz.

Artık--tahmin grafiklerinde ayrık bantlar ve yaklaşık $-4{,}24$'e
uzanan negatif artıklar vardır. Sıfır notlu gözlemlerde
artık $0-\widehat{G3}$ olduğundan aşağı eğimli bir sıra oluşur.
Histogramlarda sıfır çevresinde yoğunlaşma ve uzun negatif kuyruk
görülür. Bu görünüm normalliği ve sabit varyansı doğrulamış sayılmaz;
bantlar tek başına belirli bir varsayım ihlalinin kanıtı değildir.
Yazılımların histogram sınıf sınırları aynı olmayabilir; sütunların
farklı görünmesi sayısal sonuçların çeliştiği anlamına gelmez.
Aykırı görünen kayıtlar yalnız grafik nedeniyle silinmez.

\Needspace{12\baselineskip}
\begin{outputguidebox}
Tablo~\ref{tab:b13-dogrulanmis-katsayi} ile kurulan model
\[
\widehat{G3}=-1{,}652804+1{,}106256\,G1
\]
biçimindedir. G1'de bir puan daha yüksek değer, modelde ortalama
G3'ün yaklaşık 1,106 puan daha yüksek olmasıyla ilişkilidir;
G1'i bir puan artırmanın G3'ü bu miktarda artıracağı söylenmez.
G1'in gözlenen aralığı 3--19'dur; sabit terim G1=0'a karşılık
geldiğinden gözlenen aralığın dışındaki bir kesişimdir.

Tablo~\ref{tab:b13-dogrulanmis-uyum} içindeki $R^2=0{,}642351$,
bu örneklemde G3 değişkenliğinin yaklaşık yüzde 64,24'ünün modelle
açıklanan payıdır. Öğrencilerin yüzde 64'ünün doğru tahmin edildiği
veya yeni öğrencilerde aynı başarının elde edileceği anlamına gelmez.
G1=10 için 9,409757, modelin ortalama tahminidir; bireysel tahmin
aralığı değildir. Bu değer ara hesaplarda yuvarlanmamış katsayılarla
hesaplanmıştır.

Küçük $p$ değeri model varsayımlarını doğrulamaz. Negatif artık
kuyruğu klasik test ve aralıkların sınırlılıklarıyla birlikte
raporlanmalıdır. Farklı öğrencilerin bağımsızlığı grafiklerden
anlaşılamaz; okul, önceki koşullar ve diğer karıştırıcılar için
düzeltme yapılmamıştır. Notlar aynı eğitim sürecinin ilişkili
ölçümleridir. İki okulla sınırlı bu verideki uyum, nedensel etkiyi,
temsil gücünü veya yeni bir okulda tahmin başarısını kanıtlamaz.
\end{outputguidebox}

\Needspace{10\baselineskip}
\textbf{Raporlama örneği (doğrulanmış çıktılarla).}
``395 öğrencide birinci dönem ve yıl sonu notları arasında pozitif
doğrusal ilişki bulunmuştur ($r=0{,}801468$,
$p\approx9{,}00143\times10^{-90}$). Sabit terimli basit regresyonda
G1 eğimi 1,106256'dır (\%95 güven aralığı
$[1{,}024393;\ 1{,}188119]$); $R^2=0{,}642351$ ve
$F(1,393)=705{,}842247$ elde edilmiştir. Modelin artık standart
hatası 2,743359 puandır. Sıfır yıl sonu notu bulunan 38 kayıt
korunmuştur. Artık grafiklerinde negatif kuyruk ve uzak noktalar
görülmüş; normallik ve sabit varyans doğrulanmış kabul edilmemiştir.
Klasik test ve güven aralığı bu sınırlılıklarla yorumlanmıştır.
Okul ve diğer karıştırıcılar için düzeltme yapılmamış; bulgular
nedensel etki veya yeni öğrencilerde doğrulanmış tahmin başarısı
olarak sunulmamıştır.''

\textbf{Görev.} Standartlaştırılmış Beta ile B eğimini ayırt edin.
Yuvarlanmamış sonuçlarla $r^2=R^2$ ve $t^2=F$ eşitliklerini kontrol
edin. G1=10 için ortalama tahmini ile bireysel tahmin aralığının
neden farklı olduğunu açıklayın. Artıkların negatif kuyruğunu
raporlayın; sıfır notları yalnız model uyumunu artırmak için silmeyin.
Eğitim yılı başlamadan yapılacak bir tahminde henüz oluşmamış \texttt{g1}
notunu kullanmanın neden uygun olmayacağını açıklayın.

\section{R ile çözümlü uygulama}
\label{sec:r-b13}

\textbf{Soru.} Bölümdeki 16 gözlemlik çalışma saati--son test örneğini
R'ye aktarın. Pearson korelasyonunu, regresyon doğrusunu, $R^2$'yi ve altı
saat için öngörüyü bulun; artık grafiklerini inceleyin.

\begin{lstlisting}[style=prismR]
veri <- data.frame(
  saat = c(5, 3, 7, 6, 6, 2, 8, 4, 4, 5, 3, 7, 3, 5, 2, 6),
  puan = c(60, 54, 68, 63, 65, 52, 70, 58,
           57, 61, 52, 64, 55, 60, 50, 63)
)
cor.test(veri$saat, veri$puan, method = "pearson")
model <- lm(puan ~ saat, data = veri)
summary(model)
confint(model)
yeni <- data.frame(saat = 6)
predict(model, newdata = yeni, interval = "confidence")
predict(model, newdata = yeni, interval = "prediction")
plot(puan ~ saat, data = veri)
abline(model, col = "blue")
plot(fitted(model), resid(model),
     xlab = "Uydurulan deger", ylab = "Artik")
abline(h = 0, lty = 2)
qqnorm(resid(model))
qqline(resid(model))
\end{lstlisting}

\begin{outputguidebox}
\textbf{Çözüm ve kontrol değerleri.} $r=0{,}9806$ ve
$R^2=0{,}9616$'dır. Regresyon doğrusu
$\hat y=44{,}5980+3{,}1373\,\text{saat}$ olur. Altı saat için nokta
öngörüsü 63,4216 puandır. Eğimin sıfır olduğu hipotezinin çift yönlü
\pval-değeri yaklaşık $2{,}62\times10^{-11}$'dir; bu sayı sıfır değildir.
\end{outputguidebox}

\textbf{Yorum.} \texttt{confidence} aynı saatteki evren ortalaması için,
\texttt{prediction} ise yeni bir bireyin puanı için aralık verir. İkincisi
bireysel değişkenliği de içerdiğinden daha geniştir. Doğrusallık, sabit hata
varyansı, bağımsızlık ve çıkarım için hata dağılımı değerlendirilmelidir.

\textbf{Pekiştirme ve yanıt.} Altı saat gözlenen 2--8 saat aralığındadır;
20 saat için öngörü dışa taşırmadır. Güçlü korelasyon, ek bir saat çalışmanın
her öğrencinin puanını nedensel olarak 3,1373 artırdığını göstermez.

\section{Bölüm özeti}

\begin{itemize}
  \item Saçılım grafiği ilişki analizinin ilk ve vazgeçilmez adımıdır.
  \item Korelasyon doğrusal ilişkinin yönünü ve gücünü, regresyon eğimi ise birimli ortalama değişimi gösterir.
  \item Artıklar gözlenen değerlerle model öngörüleri arasındaki farklardır.
  \item Güçlü ilişki veya yüksek $R^2$, tek başına nedensellik ve dış örneklem başarısı göstermez.
  \item Eğrilik, huni örüntüsü ve etkili gözlemler katsayıların ve klasik
  standart hataların yorumunu değiştirebilir.
  \item Ortalama yanıt aralığı ile bireysel öngörü aralığı farklı hedeflere
  yönelir; bireysel öngörü aralığı daha geniştir.
\end{itemize}

\section*{Alıştırmalar}

\begin{enumerate}
  \item Korelasyonun neden nedensellik göstermediğini açıklayın.
  \item Regresyon eğimini birimiyle yorumlayın.
  \item Bir teknoloji destekli analiz raporunda bulunması gereken öğeleri yazın.
  \item $R^2=0{,}962$ değerini doğru biçimde yorumlayın.
  \item Regresyon doğrusunu gözlenen veri aralığının dışında kullanmanın riskini açıklayın.
  \item Altı saat için hesaplanan pozitif artığın anlamını açıklayın.
  \item Eğim güven aralığı $[2{,}78;3{,}50]$ değerini çalışma süresi bağlamında yorumlayın.
  \item Eğrisel bir ilişkide Pearson korelasyonunun neden yanıltıcı olabileceğini tartışın.
  \item Ortalama yanıt güven aralığı ile bireysel öngörü aralığını karşılaştırın.
  \item $r=-0{,}60$ değerinin yönünü ve $R^2$ karşılığını yorumlayın.
  \item $r=0{,}50$, $s_X=4$ ve $s_Y=12$ ise regresyon eğimini hesaplayın.
  \item $\bar x=10$, $\bar y=40$ ve $b_1=2$ için sabit terimi ve regresyon
  doğrusunu bulun.
  \item $\hat Y=30+4X$ modelinde $X=6$ ve gözlenen $Y=57$ için öngörü ve
  artığı hesaplayın.
  \item $n=30$ ve $r=0{,}35$ için korelasyon test istatistiğini hesaplayın.
  \item Artık grafiğindeki U biçimi ile huni biçiminin işaret ettiği sorunları
  karşılaştırın.
  \item Yüksek kaldıraçlı bir gözlemin neden mutlaka etkili veya hatalı
  olmayabileceğini açıklayın.
  \item $R^2=0{,}08$ ve çok küçük \pval-değeri bulunan büyük bir örneklemi
  istatistiksel kanıt ve öngörü başarısı bakımından yorumlayın.
  \item Çalışma süresi 2--8 saat aralığında gözlenmiş bir modelle 20 saat için
  öngörü yapmanın sakıncasını açıklayın.
\end{enumerate}